\ifdefined\gkver\else\def\gkver{student}\fi
\documentclass[\gkver]{gaokaozhenti}
\begin{document}

\chapter{2025年河南}

\begin{enumerate}
\item
%% number: 1
%% paperName: 2025年河南高考第 1 题 4 分
%% typeId: 1
%% type: 单选题
%% chapter: 运动和力
%% point: 牛顿第二定律
%% method: 无
%% score: 4
%% degree: 750
%% duplicateId: 0
%% body:
野外高空作业时，使用无人机给工人运送零件。如图，某次运送过程中的一段时间内，无人机向左水平飞行，零件用轻绳悬挂于无人机下方，并相对于无人机静止，轻绳与竖直方向成一定角度。忽略零件所受空气阻力，则在该段时间内\xzanswer{D}

\onepicture[w=4.20cm]{\includesvg{figs/01}}

\fourchoices[answer=D]
{无人机做匀速运动}
{零件所受合外力为零}
{零件的惯性逐渐变大}
{零件的重力势能保持不变}

%% answer: D
%% memo:
\memoanswer{
D．无人机沿水平方向飞行，零件相对于无人机静止，也沿水平方向飞行做直线运动，故零件的高度不变，可知零件的重力势能保持不变，D 正确；

AB．对零件受力分析，受重力和绳子的拉力，由于零件沿水平方向做直线运动，可知合外力沿水平方向，提供水平方向的加速度。零件水平向左做匀加速直线运动，AB 错误；

C．惯性的大小只与质量有关，零件的质量不变，故零件的惯性不变，C 错误。

故选 D。
}

\item
%% number: 2
%% paperName: 2025年河南高考第 2 题 4 分
%% typeId: 1
%% type: 单选题
%% chapter: 光学
%% point: 光的折射
%% method: 无
%% score: 4
%% degree: 650
%% duplicateId: 0
%% body:
折射率为 $\sqrt{2}$ 的玻璃圆柱水平放置，平行于其横截面的一束光线从顶点入射，光线与竖直方向的夹角为 $45^\circ$，如图所示。该光线从圆柱内射出时，与竖直方向的夹角为（不考虑光线在圆柱内的反射）\xzanswer{B}

\onepicture[w=3.00cm]{\includesvg{figs/02}}

\fourchoices[answer=B]
{$0^\circ$}
{$15^\circ$}
{$30^\circ$}
{$45^\circ$}

%% answer: B
%% memo:
\memoanswer{
设光线射入圆柱体时的折射角为 $\theta$，根据光的折射定律可知 $n = \frac{\sin 45^\circ}{\sin \theta}$

解得 $\theta = 30^\circ$

\onepicture[w=3.00cm]{figs/02_an}

如图，根据几何关系可知光线射出圆柱体时的入射角 $i = \theta = 30^\circ$

则法线与竖直方向的夹角 $\alpha = \theta + i = 60^\circ$

根据光的折射定律可知 $n = \frac{\sin r}{\sin i}$

解得光线射出圆柱体时的折射角 $r = 45^\circ$

光线从圆柱体内射出时，与竖直方向的夹角为 $\beta = \alpha - r = 15^\circ$

故选 B。
}

\item
%% number: 3
%% paperName: 2025年河南高考第 3 题 4 分
%% typeId: 1
%% type: 单选题
%% chapter: 曲线运动
%% point: 天体运动
%% method: 无
%% score: 4
%% degree: 550
%% duplicateId: 0
%% body:
2024 年天文学家报道了他们新发现的一颗类地行星 Gliese12b，它绕其母恒星的运动可视为匀速圆周运动。已知 Gliese12b 轨道半径约为日地距离的 $\frac{1}{14}$，其母恒星质量约为太阳质量的 $\frac{2}{7}$，则 Gliese12b 绕其母恒星的运动周期约为\xzanswer{A}

\fourchoices[answer=A]
{13 天}
{27 天}
{64 天}
{128 天}

%% answer: A
%% memo:
\memoanswer{
地球绕太阳运行的周期约为 365 天，根据万有引力提供向心力得 $\frac{GM_{0}m}{r_{0}^{2}} = m\frac{4\pi^{2}}{T_{0}^{2}}r_{0}$

已知 $r = \frac{1}{14}r_{0}$，$M = \frac{2}{7}M_{0}$，同理得 $\frac{GMm}{r^{2}} = m\frac{4\pi^{2}}{T^{2}}r$

整理得 $\frac{T^{2}}{T_{0}^{2}} = \frac{r^{3}M_{0}}{r_{0}^{3}M}$

代入数据得 $T = \frac{1}{28}T_{0} \approx 13$ 天

故选 A。
}

\item
%% number: 4
%% paperName: 2025年河南高考第 4 题 4 分
%% typeId: 1
%% type: 单选题
%% chapter: 电场
%% point: 等势面
%% method: 无
%% score: 4
%% degree: 550
%% duplicateId: 0
%% body:
如图，在与纸面平行的匀强电场中有 a、b、c 三点，其电势分别为 $6 \UV$、$4 \UV$、$2 \UV$；a、b、c 分别位于纸面内一等边三角形的顶点上。下列图中箭头表示 a 点电场的方向，则正确的是\xzanswer{C}

\onepicture[w=2.60cm]{\includesvg{figs/04a}}

\fourchoices[answer=C, ispicture=true, h=2.20cm]
{figs/04b.png}
{figs/04c.png}
{figs/04d.png}
{figs/04e.png}

%% answer: C
%% memo:
\memoanswer{
匀强电场中任意两点间的中点电势等于这两点的平均值，可知 ac 中点 d 的电势与 b 点相同，bd 的连线为该匀强电场的等势面。电场线垂直于等势面且由高电势指向低电势，故电场线沿 ac 方向且由 a 指向 c，C 选项正确。

\onepicture[w=2.60cm]{figs/04_an}

故选 C。
}

\item
%% number: 5
%% paperName: 2025年河南高考第 5 题 4 分
%% typeId: 1
%% type: 单选题
%% chapter: 电磁感应
%% point: 自感与互感，涡流
%% method: 无
%% score: 4
%% degree: 650
%% duplicateId: 0
%% body:
如图，一金属薄片在力 $F$ 作用下自左向右从两磁极之间通过。当金属薄片中心运动到 N 极的正下方时，沿 N 极到 S 极的方向看，下列图中能够正确描述金属薄片内涡电流绕行方向的是\xzanswer{C}

\onepicture[w=3.60cm]{figs/05a.png}

\fourchoices[answer=C, ispicture=true, h=1.40cm]
{figs/05b.png}
{figs/05c.png}
{figs/05d.png}
{figs/05e.png}

%% answer: C
%% memo:
\memoanswer{
根据题意当金属薄片中心运动到 N 极正下方时，薄片右侧的磁通量在减小，左侧磁通量在增加，由于两极间的磁场竖直向下，根据楞次定律可知此时薄片右侧的涡电流方向为顺时针，薄片左侧的涡电流方向为逆时针。

故选 C。
}

\item
%% number: 6
%% paperName: 2025年河南高考第 6 题 4 分
%% typeId: 1
%% type: 单选题
%% chapter: 原子和原子核
%% point: 半衰期
%% method: 无
%% score: 4
%% degree: 550
%% duplicateId: 0
%% body:
由于宇宙射线的作用，在地球大气层产生有铍的两种放射性同位素 \ce{^{7}_{4}Be} 和 \ce{^{10}_{4}Be}。测定不同高度大气中单位体积内二者的原子个数比，可以研究大气环境的变化。已知 \ce{^{7}_{4}Be} 和 \ce{^{10}_{4}Be} 的半衰期分别约为 53 天和 139 万年。在大气层某高度采集的样品中，研究人员发现 \ce{^{7}_{4}Be} 和 \ce{^{10}_{4}Be} 的总原子个数经过 106 天后变为原来的 $\frac{3}{4}$，则采集时该高度的大气中 \ce{^{7}_{4}Be} 和 \ce{^{10}_{4}Be} 的原子个数比约为\xzanswer{B}

\fourchoices[answer=B]
{$1:4$}
{$1:2$}
{$3:4$}
{$1:1$}

%% answer: B
%% memo:
\memoanswer{
设采集时大气中有 $x$ 个 \ce{^{7}_{4}Be} 原子和 $y$ 个 \ce{^{10}_{4}Be} 原子，由于 \ce{^{10}_{4}Be} 的半衰期为 139 万年，故经过 106 天后 \ce{^{10}_{4}Be} 原子的衰变个数可以忽略不计，\ce{^{7}_{4}Be} 的半衰期为 53 天，故经过 106 天后剩余数量为 $x \cdot \left(\frac{1}{2}\right)^{2}$，故可得 $\frac{x \cdot \left(\frac{1}{2}\right)^{2} + y}{x + y} = \frac{3}{4}$

解得 $\frac{x}{y} = \frac{1}{2}$

故选 B。
}

\item
%% number: 7
%% paperName: 2025年河南高考第 7 题 4 分
%% typeId: 1
%% type: 单选题
%% chapter: 运动和力
%% point: 动量守恒定律
%% method: 无
%% score: 4
%% degree: 450
%% duplicateId: 0
%% body:
两小车 P、Q 的质量分别为 $m_{P}$ 和 $m_{Q}$，将它们分别与小车 N 沿直线做碰撞实验，碰撞前后的速度 $v$ 随时间 $t$ 的变化分别如图~\ref{2025河南07a} 和图~\ref{2025河南07b} 所示。小车 N 的质量为 $m_{N}$，碰撞时间极短，则\xzanswer{D}

\twopicture[labelA=2025河南07a, labelB=2025河南07b, numA=1, numB=2, w=5.00cm]{figs/07a.png}{figs/07b.png}

\fourchoices[answer=D]
{$m_{P} > m_{N} > m_{Q}$}
{$m_{N} > m_{P} > m_{Q}$}
{$m_{Q} > m_{P} > m_{N}$}
{$m_{Q} > m_{N} > m_{P}$}

%% answer: D
%% memo:
\memoanswer{
PN 碰撞时，根据碰撞前后动量守恒有 $m_{\mathrm{P}}v_{\mathrm{P}} + m_{\mathrm{N}}v_{\mathrm{N}} = m_{\mathrm{P}}v_{\mathrm{P}}' + m_{\mathrm{N}}v_{\mathrm{N}}'$

即 $m_{\mathrm{P}}\left(v_{\mathrm{P}} - v_{\mathrm{P}}'\right) = m_{\mathrm{N}}\left(v_{\mathrm{N}}' - v_{\mathrm{N}}\right)$

根据图像可知 $\left(v_{\mathrm{P}} - v_{\mathrm{P}}'\right) > \left(v_{\mathrm{N}}' - v_{\mathrm{N}}\right)$，故 $m_{\mathrm{N}} > m_{\mathrm{P}}$；

同理，QN 碰撞时，根据碰撞前后动量守恒有 $m_{Q}v_{Q} + m_{N}v_{N} = m_{Q}v_{Q}' + m_{N}v_{N}'$

即 $m_{Q}\left(v_{Q} - v_{Q}'\right) = m_{N}\left(v_{N}' - v_{N}\right)$

根据图像可知 $\left(v_{\mathrm{Q}} - v_{\mathrm{Q}}'\right) < \left(v_{\mathrm{N}}' - v_{\mathrm{N}}\right)$，故 $m_{Q} > m_{N}$；

故 $m_{Q} > m_{N} > m_{P}$

故选 D。
}

\item
%% number: 8
%% paperName: 2025年河南高考第 8 题 6 分
%% typeId: 2
%% type: 多选题
%% chapter: 机械波
%% point: 波长频率波速
%% method: 无
%% score: 6
%% degree: 550
%% duplicateId: 0
%% body:
贾湖骨笛是河南博物院镇馆之宝之一，被誉为“中华第一笛”。其中一支骨笛可以发出 $A_{5}$、$B_{5}$、$C_{6}$、$D_{6}$、$E_{6}$ 等音。已知 $A_{5}$ 音和 $D_{6}$ 音所对应的频率分别为 $880 \UHz$ 和 $1175 \UHz$，则\xzanswer{AD}

\fourchoices[answer=AD]
{在空气中传播时，$A_{5}$ 音的波长大于 $D_{6}$ 音的}
{在空气中传播时，$A_{5}$ 音的波速小于 $D_{6}$ 音的}
{由空气进入水中，$A_{5}$ 音和 $D_{6}$ 音的频率都变大}
{由空气进入水中，$A_{5}$ 音的波长改变量大于 $D_{6}$ 音的}

%% answer: AD
%% memo:
\memoanswer{
B．声音在相同介质中的传播速度相同，因此 $A_{5}$ 和 $D_{6}$ 的传播速度相同，B 错误；

A．由 $\lambda = \frac{v}{f}$ 可知，$A_{5}$ 的波长大于 $D_{6}$ 的波长，A 正确；

C．由空气进入水中，频率不发生变化，C 错误；

D．空气中 $\lambda_{0} = \frac{v}{f}$，在水中 $\lambda = \frac{v'}{f}$，其中声音的速度只与介质有关，即在水中它们的速度大小也一样，则可得到波长的改变量为 $\Delta\lambda = \frac{v' - v}{f}$。可知频率越小其对应的波长改变量越大，D 正确。

故选 AD。
}

\item
%% number: 9
%% paperName: 2025年河南高考第 9 题 6 分
%% typeId: 2
%% type: 多选题
%% chapter: 电磁感应
%% point: 右手定则
%% method: 无
%% score: 6
%% degree: 450
%% duplicateId: 0
%% body:
手机拍照时手的抖动产生的微小加速度会影响拍照质量，光学防抖技术可以消除这种影响。如图，镜头仅通过左、下两侧的弹簧与手机框架相连，两个相同线圈 c、d 分别固定在镜头右、上两侧，c、d 中的一部分处在相同的匀强磁场中，磁场方向垂直纸面向里。拍照时，手机可实时检测手机框架的微小加速度 $a$ 的大小和方向，依此自动调节 c、d 中通入的电流 $I_{c}$ 和 $I_{d}$ 的大小和方向（无抖动时 $I_{c}$ 和 $I_{d}$ 均为零），使镜头处于零加速度状态。下列说法正确的是\xzanswer{BC}

\onepicture[h=4.50cm]{figs/09.png}

\fourchoices[answer=BC]
{若 $I_{c}$ 沿顺时针方向，$I_{d} = 0$，则表明 $a$ 的方向向右}
{若 $I_{d}$ 沿顺时针方向，$I_{c} = 0$，则表明 $a$ 的方向向下}
{若 $a$ 的方向沿左偏上 $30^\circ$，则 $I_{c}$ 沿顺时针方向，$I_{d}$ 沿逆时针方向且 $I_{c} > I_{d}$}
{若 $a$ 的方向沿右偏上 $30^\circ$，则 $I_{c}$ 沿顺时针方向，$I_{d}$ 沿顺时针方向且 $I_{c} < I_{d}$}

%% answer: BC
%% memo:
\memoanswer{
A．$I_{c}$ 顺时针而 $I_{d} = 0$，则 c 线圈受到向右的安培力，故手机的加速度是向左，使镜头处于零加速度状态，A 错误；

B．$I_{d}$ 顺时针而 $I_{c} = 0$，则 d 线圈受到向上的安培力，镜头处于零加速度状态，则手机加速度方向向下，B 正确；

C．若 $a$ 的方向左偏上 $30^\circ$，说明手机框架给镜头向上以及向左的作用力，要使得镜头处于零加速度状态，线圈 c 需要受到向右的安培力 $F_{c}$、线圈 d 需要受到向下的安培力 $F_{d}$，且 $F_{c} > F_{d}$，故可知 $I_{c}$ 顺时针 $I_{d}$ 逆时针，由 $F = BIl$ 可知 $I_{c} > I_{d}$，C 正确；

D．若 $a$ 的方向右偏上 $30^\circ$，说明手机框架给镜头向上以及向右的作用力，且向右的分力大于向上的分力要使得镜头处于零加速度状态，线圈 c 需要受到向左的安培力 $F_{c}$、线圈 d 需要受到向下的安培力 $F_{d}$，且 $F_{c} > F_{d}$，可知 $I_{c}$ 逆时针 $I_{d}$ 逆时针，且 $I_{c} > I_{d}$，D 错误。

故选 BC。
}

\item
%% number: 10
%% paperName: 2025年河南高考第 10 题 6 分
%% typeId: 2
%% type: 多选题
%% chapter: 气体性质
%% point: 液柱移动定性分析
%% method: 无
%% score: 6
%% degree: 350
%% duplicateId: 0
%% body:
如图，一圆柱形汽缸水平固置，其内部被活塞 M、P、N 密封成两部分，活塞 P 与汽缸壁均绝热且两者间无摩擦。平衡时，P 左、右两侧理想气体的温度分别为 $T_{1}$ 和 $T_{2}$，体积分别为 $V_{1}$ 和 $V_{2}$，$T_{1} < T_{2}$，$V_{1} < V_{2}$。则\xzanswer{AC}

\onepicture[h=1.80cm]{figs/10.png}

\fourchoices[answer=AC]
{固定 M、N，若两侧气体同时缓慢升高相同温度，P 将右移}
{固定 M、N，若两侧气体同时缓慢升高相同温度，P 将左移}
{保持 $T_{1}$、$T_{2}$ 不变，若 M、N 同时缓慢向中间移动相同距离，P 将右移}
{保持 $T_{1}$、$T_{2}$ 不变，若 M、N 同时缓慢向中间移动相同距离，P 将左移}

%% answer: AC
%% memo:
\memoanswer{
AB．由题干可知初始左右气体的压强相同，假设在升温的过程中 P 板不发生移动，则由等容过程 $\frac{p}{T} = \frac{\Delta p}{\Delta T} \Rightarrow \Delta p = \frac{p}{T} \Delta T$

可得左侧气体压强增加量多，则 P 板向右移动；A 正确 B 错误；

CD．保持温度不变移动相同的距离时

$\frac{pV_{1}}{T_{1}} = C_{1}$，$p = \frac{C_{1}T_{1}}{V_{1}}$，$\frac{C_{1}T_{1}}{V_{1}} = \frac{C_{2}T_{2}}{V_{2}}$，$\frac{V_{1}}{C_{1}T_{1}} = \frac{V_{2}}{C_{2}T_{2}}$

同理得，$\frac{V_{1} - \Delta V}{C_{1}T_{1}} < \frac{V_{2} - \Delta V}{C_{2}T_{2}}$，$\frac{C_{1}T_{1}}{V_{1} - \Delta V} > \frac{C_{2}T_{2}}{V_{2} - \Delta V}$

若 P 不移动，则 $p_{1} = \frac{C_{1}T_{1}}{V_{1} - \Delta V}$、$p_{2} = \frac{C_{2}T_{2}}{V_{2} - \Delta V}$，故 $p_{1} > p_{2}$，则 P 将向右移动，C 正确 D 错误。

故选 AC。
}

\item
%% number: 11
%% paperName: 2025年河南高考第 11 题 6 分
%% typeId: 4
%% type: 实验题
%% chapter: 电路
%% point: 传感器
%% method: 无
%% score: 6
%% degree: 550
%% duplicateId: 0
%% body:
实验小组研究某热敏电阻的特性，并依此利用电磁铁、电阻箱等器材组装保温箱。该热敏电阻阻值随温度的变化曲线如图~\ref{2025河南11a} 所示，保温箱原理图如图~\ref{2025河南11b} 所示。回答下列问题：

\begin{enumerate}
	\item
	图~\ref{2025河南11a} 中热敏电阻的阻值随温度的变化关系是\tkanswer{非线性}（填“线性”或“非线性”）的。

	\item
	存在一个电流值 $I_{0}$，若电磁铁线圈的电流小于 $I_{0}$，衔铁与上固定触头 a 接触；若电流大于 $I_{0}$，衔铁与下固定触头 b 接触。保温箱温度达到设定值后，电磁铁线圈的电流在 $I_{0}$ 附近上下波动，加热电路持续地断开、闭合，使保温箱温度维持在设定值。则图~\ref{2025河南11b} 中加热电阻丝的 c 端应该与触头\tkanswer{a}（填“a”或“b”）相连接。

	\item
	当保温箱的温度设定在 $50\UdC$ 时，电阻箱旋钮的位置如图~\ref{2025河南11c} 所示，则电阻箱接入电路的阻值为\tkanswer{130.0} $\UO$。

	\item
	若要把保温箱的温度设定在 $100\UdC$，则电阻箱接入电路的阻值应为\tkanswer{210.0} $\UO$。
\end{enumerate}

\threepicture[labelA=2025河南11a, labelB=2025河南11b, labelC=2025河南11c, numA=1, numB=2, numC=3, wA=4.50cm, wB=5.50cm, wC=3.00cm, align=right]{figs/11a.png}{figs/11b.png}{figs/11c.png}

%% answer:
\jdanswer{
\begin{enumerate}
	\item
	非线性

	\item
	a

	\item
	130.0

	\item
	210.0
\end{enumerate}
}
%% memo:
\memoanswer{
（1）根据图 1 可知热敏电阻的阻值随温度的变化关系是非线性的。

（2）根据图 1 可知温度升高，热敏电阻的阻值变小，根据欧姆定律可知流过电磁铁线圈的电流变大，衔铁与下固定触头 b 接触，此时加热电阻丝电路部分断开连接，停止加热，可知图 2 中加热电阻丝的 c 端应该与触头 a 相连接。

（3）由图 3 可知电阻箱接入电路的阻值为 $100 \times 1 \UO + 10 \times 3 \UO = 130.0 \UO$。

（4）根据（3）可知，当温度为 $50\UdC$ 时，热敏电阻的阻值为 $180 \UO$，电阻箱接入的电阻为 $130 \UO$，当温度为 $100\UdC$ 时，热敏电阻的阻值为 $100 \UO$，要使得电流值 $I_{0}$ 不变，则在电流为 $I_{0}$ 时，控制电路的总电阻不变，则此时电阻箱的电阻为 $180 \UO + 130 \UO - 100 \UO = 210 \UO$。
}

\item
%% number: 12
%% paperName: 2025年河南高考第 12 题 9 分
%% typeId: 4
%% type: 实验题
%% chapter: 功和能
%% point: DIS研究机械能守恒定律
%% method: 无
%% score: 9
%% degree: 450
%% duplicateId: 0
%% body:
实验小组利用图~\ref{2025河南12a} 所示装置验证机械能守恒定律。可选用的器材有：交流电源（频率 $50 \UHz$）、铁架台、电子天平、重锤、打点计时器、纸带、刻度尺等。

\begin{enumerate}
	\item
	下列所给实验步骤中，有 4 个是完成实验必需且正确的，把它们选择出来并按实验顺序排列：\tkanswer{④①⑥⑤}（填步骤前面的序号）

	\begin{steps}
		\item
		先接通电源，打点计时器开始打点，然后再释放纸带
		\item
		先释放纸带，然后再接通电源，打点计时器开始打点
		\item
		用电子天平称量重锤的质量
		\item
		将纸带下端固定在重锤上，穿过打点计时器的限位孔，用手捏住纸带上端
		\item
		在纸带上选取一段，用刻度尺测量该段内各点到起点的距离，记录分析数据
		\item
		关闭电源，取下纸带
	\end{steps}

	\item
	图~\ref{2025河南12b} 所示是纸带上连续打出的五个点 A、B、C、D、E 到起点的距离。则打出 B 点时重锤下落的速度大小为\tkanswer{1.79} $\Ums$（保留 3 位有效数字）。

	\item
	纸带上各点与起点间的距离即为重锤下落高度 $h$，计算相应的重锤下落速度 $v$，并绘制图~\ref{2025河南12c} 所示的 $v^{2}-h$ 关系图像。理论上，若机械能守恒，图中直线应\tkanswer{通过}（填“通过”或“不通过”）原点且斜率为\tkanswer{$2g$}（用重力加速度大小 $g$ 表示）。由图~\ref{2025河南12c} 得直线的斜率 $k = $\tkanswer{19.1}（保留 3 位有效数字）。

	\item
	定义单次测量的相对误差 $\eta = \left| \frac{E_{\mathrm{p}} - E_{\mathrm{k}}}{E_{\mathrm{p}}} \right| \times 100\%$，其中 $E_{p}$ 是重锤重力势能的减小量，$E_{k}$ 是其动能增加量，则实验相对误差为 $\eta = $\tkanswer{$\frac{2g - k}{2g}$} $\times 100\%$（用字母 $k$ 和 $g$ 表示）；当地重力加速度大小取 $g = 9.80 \Umsq$，则 $\eta = $\tkanswer{3.1} $\%$（保留 2 位有效数字），若 $\eta < 5\%$，可认为在实验误差允许的范围内机械能守恒。
\end{enumerate}

\threepicture[labelA=2025河南12a, labelB=2025河南12b, labelC=2025河南12c, numA=1, numB=2, numC=3, wA=2.80cm, wB=6.50cm, wC=4.00cm, align=right]{figs/12a.png}{figs/12b.png}{figs/12c.png}

%% answer:
\jdanswer{
\begin{enumerate}
	\item
	④①⑥⑤

	\item
	1.79

	\item
	通过，$2g$，19.1

	\item
	$\frac{2g - k}{2g}$，3.1
\end{enumerate}
}
%% memo:
\memoanswer{
（1）实验步骤为：将纸带下端固定在重锤上，穿过打点计时器的限位孔，用手捏住纸带上端，先接通电源，打点计时器开始打点，然后再释放纸带，关闭电源，取下纸带，在纸带上选取一段，用刻度尺测量该段内各点到起点的距离，记录分析数据，根据 $mgh = \frac{1}{2}mv^{2}$ 可知质量可以约掉，不需要用电子天平称量重锤的质量。

故选择正确且正确排序为④①⑥⑤这四步。

（2）根据题意可知纸带上相邻计数点时间间隔 $T = \frac{1}{f} = 0.02 \Us$

根据匀变速直线运动中间时刻瞬时速度等于该过程平均速度可得 $v_{B} = \frac{h_{AC}}{2T}$

代入数据可得 $v_{B} \approx 1.79 \Ums$

（3）根据 $mgh = \frac{1}{2}mv^{2}$

整理可得 $v^{2} = 2g \cdot h$

可知理论上，若机械能守恒，图中直线应通过原点，且斜率 $k = 2g$

由图 3 得直线的斜率 $k = \frac{5.6 - 1.5}{0.295 - 0.08} \approx 19.1$

（4）根据题意有 $\eta = \frac{mgh - \frac{1}{2}mv^{2}}{mgh} \times 100\%$

可得 $\eta = \frac{2g - k}{2g} \times 100\%$

当地重力加速度大小取 $g = 9.80 \Umsq$，代入数据可得 $\eta \approx 3.1\%$。
}

\item
%% number: 13
%% paperName: 2025年河南高考第 13 题 10 分
%% typeId: 6
%% type: 计算题
%% chapter: 电场
%% point: 带电粒子的偏转
%% method: 无
%% score: 10
%% degree: 550
%% duplicateId: 0
%% body:
流式细胞仪可对不同类型的细胞进行分类收集，其原理如图所示。仅含有一个 A 细胞或 B 细胞的小液滴从喷嘴喷出（另有一些液滴不含细胞），液滴质量均为 $m = 2.0 \times 10^{-10} \Ukg$。当液滴穿过激光束、充电环时被分类充电，使含 A、B 细胞的液滴分别带上正、负电荷，电荷量均为 $q = 1.0 \times 10^{-13} \UC$。随后，液滴以 $v = 2.0 \Ums$ 的速度竖直进入长度为 $l = 2.0 \times 10^{-2} \Um$ 的电极板间，板间电场均匀、方向水平向右，电场强度大小为 $E = 2.0 \times 10^{5} \UNC$。含细胞的液滴最终被分别收集在极板下方 $h = 0.1 \Um$ 处的 A、B 收集管中。不计重力、空气阻力以及带电液滴间的作用。求：

\begin{enumerate}
	\item
	含 A 细胞的液滴离开电场时偏转的距离；
	\item
	A、B 细胞收集管的间距。
\end{enumerate}

\onepicture[h=5.50cm, align=right]{figs/13.png}

%% answer:
\jdanswer{
\begin{enumerate}
	\item
	$5 \times 10^{-3} \Um$

	\item
	$0.11 \Um$
\end{enumerate}
}
%% memo:
\memoanswer{
（1）由题意可知含 A 细胞的液滴在电场中做类平抛运动，垂直于电极板方向则 $l = vt_{1}$，沿电极板方向 $x_{1} = \frac{1}{2}at_{1}^{2}$

由牛顿第二定律 $qE = ma$

解得含 A 细胞的液滴离开电场时偏转的距离为 $x_{1} = 5 \times 10^{-3} \Um$

（2）含 A 细胞的液滴离开电场后做匀速直线运动，则 $h = vt_{2}$

则 $x_{2} = at_{1}t_{2}$

联立解得 $x_{2} = 0.05 \Um$

由对称性可知则 A、B 细胞收集管的间距 $\Delta x = 2(x_{1} + x_{2}) = 2 \times (0.005 + 0.05) \Um = 0.11 \Um$
}

\item
%% number: 14
%% paperName: 2025年河南高考第 14 题 12 分
%% typeId: 6
%% type: 计算题
%% chapter: 运动和力
%% point: 动量和能量
%% method: 无
%% score: 12
%% degree: 450
%% duplicateId: 0
%% body:
如图，在一段水平光滑直道上每间隔 $l_{1} = 3 \Um$ 铺设有宽度为 $l_{2} = 2.4 \Um$ 的防滑带。在最左端防滑带的左边缘静止有质量为 $m_{1} = 2 \Ukg$ 的小物块 P，另一质量为 $m_{2} = 4 \Ukg$ 的小物块 Q 以 $v_{0} = 7 \Ums$ 的速度向右运动并与 P 发生正碰，且碰撞时间极短。已知碰撞后瞬间 P 的速度大小为 $v = 7 \Ums$，P、Q 与防滑带间的动摩擦因数均为 $\mu = 0.5$，重力加速度大小 $g = 10 \Umsq$。求：

\begin{enumerate}
	\item
	该碰撞过程中损失的机械能；
	\item
	P 从开始运动到静止经历的时间。
\end{enumerate}

\onepicture[h=1.80cm, align=right]{figs/14.png}

%% answer:
\jdanswer{
\begin{enumerate}
	\item
	$24.5 \UJ$

	\item
	$5 \Us$
\end{enumerate}
}
%% memo:
\memoanswer{
（1）P、Q 发生正碰，由动量守恒定律 $m_{2}v_{0} = m_{2}v_{Q} + m_{1}v$

由能量守恒定律 $\frac{1}{2}m_{2}v_{0}^{2} = \frac{1}{2}m_{2}v_{Q}^{2} + \frac{1}{2}m_{1}v^{2} + \Delta E$

联立可得 $v_{Q} = 3.5 \Ums$，$\Delta E = 24.5 \UJ$

（2）对物块 P 受力分析由牛顿第二定律 $\mu m_{1}g = m_{1}a$

物块 P 在第一个防滑带上运动时，由运动学公式 $v^{2} - v_{P1}^{2} = 2al_{2}$，$v_{P1} = v - at_{1}$

解得 $v_{P1} = 5 \Ums$

则物块 P 在第一个防滑带上运动的时间为 $t_{1} = 0.4 \Us$

物块 P 在光滑的直道上做匀速直线运动，则 $l_{1} = v_{P1}t_{2}$

解得 $t_{2} = 0.6 \Us$

物块 P 在第二个防滑带上运动时，由运动学公式 $v_{P1}^{2} - v_{P2}^{2} = 2al_{2}$，$v_{P2} = v_{P1} - at_{3}$

解得 $v_{P2} = 1 \Ums$

则物块 P 在第二个防滑带上运动的时间为 $t_{3} = 0.8 \Us$

物块 P 在光滑的直道上做匀速直线运动，则 $l_{1} = v_{P2}t_{4}$

解得 $t_{4} = 3 \Us$

由以上条件可知，物块 P 最终停在第三个防滑带上，由运动学公式 $0 = v_{P2} - at_{5}$

可得物块 P 在第三个防滑带上运动的时间为 $t_{5} = 0.2 \Us$

故物块 P 从开始运动到静止经历的时间为 $t = t_{1} + t_{2} + t_{3} + t_{4} + t_{5} = 5 \Us$
}

\item
%% number: 15
%% paperName: 2025年河南高考第 15 题 17 分
%% typeId: 6
%% type: 计算题
%% chapter: 磁场
%% point: 带电粒子在组合场中的运动
%% method: 无
%% score: 17
%% degree: 350
%% duplicateId: 0
%% body:
如图，水平虚线上方区域有垂直于纸面向外的匀强磁场，下方区域有竖直向上的匀强电场。质量为 $m$、带电量为 $q$（$q > 0$）的粒子从磁场中的 a 点以速度 $v_{0}$ 向右水平发射，当粒子进入电场时其速度沿右下方向并与水平虚线的夹角为 $60^\circ$，然后粒子又射出电场重新进入磁场并通过右侧 b 点，通过 b 点时其速度方向水平向右。a、b 距水平虚线的距离均为 $h$，两点之间的距离为 $s = 3\sqrt{3}h$。不计重力。

\begin{enumerate}
	\item
	求磁感应强度的大小；
	\item
	求电场强度的大小；
	\item
	若粒子从 a 点以 $v_{0}$ 竖直向下发射，长时间来看，粒子将向左或向右漂移，求漂移速度大小。（一个周期内粒子的位移与周期的比值为漂移速度）
\end{enumerate}

\onepicture[h=3.40cm, align=right]{figs/15.png}

%% answer:
\jdanswer{
\begin{enumerate}
	\item
	$\frac{mv_{0}}{2qh}$

	\item
	$\frac{mv_{0}^{2}}{2qh}$

	\item
	$\frac{3\sqrt{3}}{6\sqrt{3} + 8\pi}v_{0}$
\end{enumerate}
}
%% memo:
\memoanswer{
（1）根据题意可知，画出粒子的运动轨迹，如图所示。

\onepicture[w=6.50cm]{figs/15_an1.png}

由题意可知 $\theta = 60^\circ$

设粒子在磁场中做圆周运动的半径为 $r$，由几何关系有 $r = r\cos\theta + h$

解得 $r = 2h$

由牛顿第二定律有 $qv_{0}B = m\frac{v_{0}^{2}}{r}$

解得 $B = \frac{mv_{0}}{2qh}$

（2）据题意，由对称性可知，粒子射出电场时，速度大小仍为 $v_{0}$，方向与水平虚线的夹角为 $60^\circ$，由几何关系可得 $AB = s - 2r\sin\theta = 3\sqrt{3}h - 2\sqrt{3}h = \sqrt{3}h$

则粒子在电场中的运动时间为 $t = \frac{AB}{v_{0}\cos\theta} = \frac{2\sqrt{3}h}{v_{0}}$

沿电场方向上，由牛顿第二定律有 $qE = ma$

由运动学公式有 $-v_{0}\sin\theta = v_{0}\sin\theta - at$

联立解得 $E = \frac{mv_{0}^{2}}{2qh}$

（3）若粒子从 a 点以 $v_{0}$ 竖直向下发射，画出粒子的运动轨迹，如图所示

\onepicture[w=6.50cm]{figs/15_an2.png}

由于粒子在磁场中运动的速度大小仍为 $v_{0}$，粒子在磁场中运动的半径仍为 $2h$，由几何关系可得，粒子进入电场时速度与虚线的夹角 $\alpha = 60^\circ$

结合小问（2）分析可知，粒子在电场中的运动时间为 $t_{1} = \frac{2\sqrt{3}h}{v_{0}}$

AB 间的距离为 $AB = \sqrt{3}h$

由几何关系可得 $BC = 2r\sin\alpha = 2\sqrt{3}h$

则 $AC = BC - AB = \sqrt{3}h$

粒子在磁场中的运动时间为 $t_{2} = \frac{360^\circ - 2\alpha}{360^\circ}\cdot\frac{2\pi r}{v_{0}} = \frac{8\pi h}{3v_{0}}$

则有 $t = t_{1} + t_{2} = \frac{\left(6\sqrt{3} + 8\pi\right)h}{3v_{0}}$

综上所述可知，粒子每隔时间 $t$ 向右移动 $\sqrt{3}h$，则漂移速度大小 $v' = \frac{\sqrt{3}h}{t} = \frac{3\sqrt{3}}{6\sqrt{3} + 8\pi}v_{0}$。
}

\end{enumerate}

\end{document}
